% Mean bibliography entries per paper, NeurIPS / ICLR / ICML.
% Data: ../../data/pgfplots/tmlr/references_per_year.csv
% Papers with no Semantic Scholar reference list are excluded rather than
% counted as zero; see src/citations/references.py.

\begin{figure}[h]
\centering
\begin{tikzpicture}
\begin{axis}[
    width=0.85\linewidth, height=5cm,
    xlabel={Year}, ylabel={Bibliography entries per paper},
    xtick distance=8,
    x tick label style={rotate=45, anchor=east, font=\tiny,
                        /pgf/number format/1000 sep={}},
    ytick align=inside,
    tick label style={font=\tiny},
    label style={font=\small},
    enlarge x limits=0.02,
    legend style={font=\tiny, at={(0.02,0.98)}, anchor=north west,
                  draw=none, fill opacity=0.7, text opacity=1},
    legend cell align=left,
  ]
  \addplot[color={rgb,255:red,44;green,160;blue,44}, dashdotted, very thick]
    table[x=year, y=neurips, col sep=comma]
    {../../data/pgfplots/tmlr/references_per_year.csv};
  \addlegendentry{NeurIPS}
  \addplot[color={rgb,255:red,31;green,119;blue,180}, densely dashed, very thick]
    table[x=year, y=iclr, col sep=comma]
    {../../data/pgfplots/tmlr/references_per_year.csv};
  \addlegendentry{ICLR}
  \addplot[color={rgb,255:red,255;green,127;blue,14}, densely dotted, very thick]
    table[x=year, y=icml, col sep=comma]
    {../../data/pgfplots/tmlr/references_per_year.csv};
  \addlegendentry{ICML}
  \draw[black, dashed, semithick] (axis cs:2022,0) -- (axis cs:2022,100);
\end{axis}
\end{tikzpicture}
\caption{\textbf{Mean bibliography length per paper at the three
  conferences.}
  \rev{Reference list length grows steadily at every venue and then
  levels off. NeurIPS rises from roughly 11 entries per paper in
  1987 to a peak of 62.9 in 2022 and holds between 59 and 61
  through 2025. ICLR, from 2013, and ICML, from 2008, converge on
  the same band, reaching 57 to 61 entries by 2025. Papers for which
  Semantic Scholar holds no reference list are excluded, and
  coverage of the most recent year is lower than of earlier years
  (73 percent of NeurIPS 2025, against 97 percent for 2022).}}
\label{fig:app-references}
\end{figure}
